\section{Related Works}
% Early seizure detection systems relied on manually extracted statistical, spectral or wavelet-based features combined with classifiers such as support vector machines, random forests and logistic regression. While these approaches often achieved good performance on smaller datasets, they depended heavily on feature engineering and struggled to generalize across patients.

% The current state of the art in seizure detection no longer relies on traditional machine learning techniques, but have rather moved on to more complex models such as transformer-based models, graph neural networks and hybrid models combining various techniques, as discussed in a survey by Xu et al \cite{XU2024128644} and the work of Garg et al \cite{review}. 
% Transformer-based models have become popular due to their self-attention feature, which means that they can learn relationships of features, even if they are far apart in the time domain, as opposed to RNN models, such as GRU or LSTM models, which have to process inputs sequentially. An expansion on this combines the transformer with a CNN, to have the CNN act as an automatic feature extractor, which will learn the most valuable features to extract from the inputs, such as spikes or sharp edges. The transformer can then learn the temporal relationships between those features. 

% Another modern approach is to use a graph neural network (GNN), which have the advantage of being able to utilise relationships between input channels, which lends itself well to multi-channel EEG readings. EEG channels exhibit complex spatial and functional relationships. GNNs can explicitly model these relationships through graph structures, allowing information to propagate between anatomically or functionally connected electrodes, as discussed by Jibon et al \cite{jibon2024sequential}. 

% Despite their strong classification performance, transformer-based and GNN-based models mostly require larger datasets, are more computationally heavy, and take longer to train. 

% Due to this, the model chosen for this paper is a hybrid CNN-LSTM model, which seeks to use the CNN as a trained feature extractor, extracting spatial features from the input, which are then passed in to the LSTM model, which learns the temporal relationships between features. This structure was chosen to strike a balance between good feature extraction and classification performance, while not being as computationally demanding as a transformer-based or GNN based model. A GNN based model is also unlikely to give much advantage, as the available dataset is only single-channel, and does not have the advantage of relationships between channels.

% You should describe the state-of-the-art (SOTA) deep neural networks relevant to your task. After presenting these methods, clearly justify the motivation behind choosing your specific algorithm. For example, if you selected MobileNet for an image classification task, explain why MobileNet was preferred over other SOTA models, highlighting its advantages in the context of your application.

% About 1 page.

Early seizure detection systems relied on manually extracted statistical, spectral, or wavelet-based features paired with traditional classifiers like support vector machines, random forests, and logistic regression. While effective on smaller datasets, these approaches depended heavily on manual feature engineering and struggled to generalize across different patients.

To overcome these limitations, the current state of the art has shifted toward deep learning architectures, notably transformer-based models, graph neural networks (GNNs), and various hybrid configurations, as reviewed by Xu et al. \cite{XU2024128644} and Garg et al. \cite{review}. Transformer-based models have gained popularity due to their self-attention mechanism, which captures long-range dependencies across the time domain. This offers a distinct advantage over recurrent neural networks (RNNs) like GRUs or LSTMs, which must process inputs sequentially and often struggle with long-term memory. To optimize this approach, recent studies have combined a transformer with varying types of convolutional networks. The works of Dong et al use a 2D-CNN \cite{DONG2025107075}, while Zhou et al use convolutional block attention \cite{Zhou2026}.

Alternatively, GNNs have emerged as a powerful tool for analyzing multi-channel EEG readings. Because EEG channels exhibit complex spatial and functional relationships, GNNs can explicitly model these interactions through graph structures, allowing information to propagate between anatomically or functionally connected electrodes, which has been explored by Jibon et al \cite{jibon2024sequential}.

Despite their high classification performance, both transformer- and GNN-based models require massive datasets, carry heavy computational overhead, and demand prolonged training times.

Consequently, this paper adopts a hybrid CNN-LSTM model. This architecture utilizes a CNN to automatically extract spatial features from the raw input, which are then passed to an LSTM to learn sequential, temporal relationships. This design strikes a practical balance between robust feature extraction and high classification accuracy without the prohibitive computational demands of transformer or GNN architectures. Furthermore, a GNN approach would yield little benefit in this study, as the available dataset consists of single-channel EEG recordings that lack the inter-channel relationships GNNs are designed to exploit.